\documentclass[11pt,a4paper]{report}
\usepackage[margin=2.5cm]{geometry}
\usepackage{amsmath,amssymb,amsfonts}
\usepackage{booktabs}
\usepackage{longtable}
\usepackage{enumitem}
\usepackage{xcolor}
\usepackage{titlesec}
\usepackage[hidelinks]{hyperref}
\usepackage{parskip}
\usepackage{fancyhdr}

\setlist{itemsep=2pt,topsep=4pt,parsep=0pt}
\definecolor{soft}{RGB}{95,99,104}
\newcommand{\change}[1]{\par\vspace{2pt}\noindent\textcolor{soft}{\small\textbf{Changed since V4.}\ #1}\par\vspace{2pt}}
\newcommand{\measured}[1]{\textbf{#1}}
\newcommand{\SOthree}{\mathrm{SO}(3)}

\pagestyle{fancy}
\fancyhf{}
\fancyhead[L]{\small\textcolor{soft}{Equivariant 3D Fracture Reassembly --- Design V5}}
\fancyhead[R]{\small\thepage}
\renewcommand{\headrulewidth}{0.2pt}

\begin{document}

\begin{titlepage}
\centering
\vspace*{4cm}
{\Huge\bfseries Design V5\par}
\vspace{0.8cm}
{\Large Architectural Framework and Shared Pipeline Design\par}
\vspace{0.4cm}
{\large Equivariant 3D Fracture Reassembly\par}
\vspace{2.5cm}
{\large M.Sc. Thesis --- Design Document\par}
\vspace{0.3cm}
{\normalsize Revision 5, superseding Design V4\par}
\vfill
\begin{minipage}{0.85\textwidth}
\small
\textit{This revision replaces Design V4. Every figure quoted here has been
measured on the implemented system or on the Breaking Bad dataset itself;
where a quantity is still an assumption, it is labelled as one. Sections that
depart from V4 carry an explicit note saying what changed and why, so that a
reader of V4 can find the deltas without re-reading the whole document.}
\end{minipage}
\vspace{2cm}
\end{titlepage}

\tableofcontents
\newpage

\chapter{Architectural Framework and Shared Pipeline Design}

\section{Introduction}
\label{sec:intro}

The objective of this research is to reconstruct fractured three-dimensional
objects from a collection of disconnected fragments. Unlike conventional shape
reconstruction, fragment reassembly requires reasoning simultaneously about
local geometry, global object structure, and relationships between fragments
that no longer share explicit graph connections.

The framework treats every fragment as an independent graph and employs an
equivariant graph neural network to learn geometric representations that remain
consistent under rotation. The central hypothesis is unchanged from V4:
reassembly should be decomposed into two problems. First, the network learns how
fragments are oriented. Second, once orientations are recovered, translations are
solved by geometric optimisation rather than direct regression. This lets the
network concentrate on geometric relationships and delegates final placement to a
mathematically grounded procedure.

The research question that motivates the whole design is one of \emph{economy}.
The current state of the art, GARF, reaches 6.1$^{\circ}$ rotation RMSE using
four H100 accelerators for 72 hours. This work asks whether equivariance can
substitute for scale: whether a network that is $\SOthree$-equivariant
\emph{by construction} needs to see far fewer examples than one that must learn
rotational invariance from data, and can therefore be trained within a two-GPU,
twelve-hour-session budget on Tesla T4 hardware.

\subsection{Scope of this revision}
\label{sub:scope}

V4 proposed two candidate backbones, Vector Neuron Graph Attention (VN-GAT) and
Tensor Field Graph Attention (TF-GAT). Only VN-GAT has been implemented and is
described here. TF-GAT remains a documented alternative that was not pursued;
nothing in the surrounding pipeline would prevent substituting it later.

Table~\ref{tab:revisions} summarises what has changed. Each row is expanded in
the section named.

\begin{table}[h]
\centering
\small
\begin{tabular}{p{4.3cm}p{5.2cm}p{5.2cm}}
\toprule
\textbf{Component} & \textbf{Design V4} & \textbf{Design V5} \\
\midrule
Inter-fragment communication & 8 virtual nodes per fragment, pooled and broadcast & Direct attention between sampled fracture-surface vertices (\S\ref{sec:tokens}, \S\ref{sub:cross}) \\
Edge features & length, midpoint, two face normals & two face normals, relative position (\S\ref{sub:edge}) \\
Fragment scale & implicit & explicit invariant gate (\S\ref{sub:scale}) \\
Connectivity & adjacency and incidence & adjacency only (\S\ref{sub:connectivity}) \\
Equivariance condition & $L(R\cdot f) = D(R)\,L(f)$ & additionally invariant to \emph{other} fragments' poses (\S\ref{sec:equivariance}) \\
Rotation representation & normalised quaternion & 6D vector pair, Gram--Schmidt (\S\ref{sub:head}) \\
Rotation loss & $\arccos$ of the trace & $\operatorname{atan2}$ formulation (\S\ref{sub:rotloss}) \\
Interface loss & $\lVert\sum_i z_i\rVert^2$ & InfoNCE over coincidence clusters (\S\ref{sub:corr}) \\
Midpoint loss & present & removed with the midpoint feature (\S\ref{sub:aux}) \\
Reference values & not specified & every term has a measured chance value (\S\ref{sec:reference}) \\
\bottomrule
\end{tabular}
\caption{Revisions from Design V4. Three of these --- the interface loss, the
rotation loss, and the equivariance condition --- are corrections of defects
rather than refinements.}
\label{tab:revisions}
\end{table}

\section{The dataset, measured}
\label{sec:dataset}

Several V4 decisions were made against assumed properties of the data. An
exhaustive pass over the corpus was therefore run before finalising this
revision: 1{,}442 objects and \measured{1{,}096{,}825 fragments}, processed in
approximately one hour on four workers.

\begin{table}[h]
\centering
\small
\begin{tabular}{lrrr}
\toprule
\textbf{Quantity} & \textbf{Median} & \textbf{Mean} & \textbf{Maximum} \\
\midrule
Vertices per fragment & 296 & 1{,}550 & 83{,}039 \\
Fracture vertices per fragment & 125 & 164 & 3{,}587 \\
Fragments per scene & 3 & 8 & 99 \\
\bottomrule
\end{tabular}
\caption{Breaking Bad, measured over the full corpus using coincidence ground
truth.}
\label{tab:dataset}
\end{table}

Dataset-wide, the fracture surface is \measured{10.6\% of vertices, 6.2\% of
faces and 7.6\% of edges}. Two reduction figures are reported throughout this
document because they answer different questions and diverge sharply: the
\emph{per-fragment mean} removes 49.9\% of vertices, weighting every fragment
equally, while the \emph{dataset aggregate} removes 89.4\%, weighting by size.
Most fragments are small, and small fragments are mostly fracture surface, so a
typical fragment loses half its vertices while the dataset loses nine tenths.

Two consequences shape the design. The heavy tail --- a maximum of 83{,}039
vertices against a median of 296 --- means any per-fragment memory budget must
be set by the tail, not the median. And because the true fracture surface is
small, the inter-fragment graph is far cheaper than V4 assumed
(\S\ref{sub:graphcost}).

\section{Geometric feature representation and graph construction}
\label{sec:features}

\subsection{Requirements}
\label{sub:requirements}

Meshes define graph structures naturally: vertices are nodes, mesh edges are
graph edges. The representation must preserve geometric information, remain
compatible with equivariant architectures, and carry information useful for
fracture matching.

The compatibility requirement is more restrictive than it appears. Vector Neuron
features are shaped $(C,3)$, and a rotation acts on the last axis as
$x \mapsto x R^{\top}$. Anything without a direction --- a length, a scalar
scale, a count --- has no legal slot in that tensor, and placing one there
silently breaks equivariance. This single constraint drives most of the changes
in this section.

\subsection{Fragment centralisation and scene standardisation}
\label{sub:centralisation}

Fragments occupy arbitrary locations. Using absolute coordinates would let the
network learn dataset placement biases rather than intrinsic geometry. For each
fragment $f_k$ with vertex set $V_k$, let
\[
c_k = \frac{1}{|V_k|}\sum_{v \in V_k} x_v, \qquad \tilde{x}_v = x_v - c_k .
\]

\change{Centralisation is now followed by a per-scene standardisation. All
fragments of a scene are divided by a single scene radius, so the network sees
shapes at a consistent scale. One divisor is used for both the ground-truth and
the perturbed copy, because a rotation does not change a radius --- which is
what allows the supervision relation of \S\ref{sub:rotloss} to survive
normalisation exactly.}

Centralisation must remain \emph{reversible}: stage two recovers per-fragment
translation, and the original centroid vectors have to be retrievable.

\subsection{Node features}
\label{sub:node}

Each node corresponds to a mesh vertex. The node representation is a
vector-valued pair
\[
V_v = [\,\tilde{x}_v,\; n_v\,] \in \mathbb{R}^{2\times 3},
\]
the centralised coordinate and the vertex normal. The coordinate describes local
shape within the fragment; the normal describes local surface orientation. This
is unchanged from V4.

\subsection{Edge features}
\label{sub:edge}

For an edge joining $u$ and $v$, V4 used
$e_{uv} = [\,l_{uv},\, m_{uv},\, n^{(1)}_{uv},\, n^{(2)}_{uv}\,]$, combining a
scalar length, a midpoint, and two face normals.

\change{The length and the midpoint are removed; a relative position is added.
The edge representation is now three vector channels,
\[
e_{uv} = [\,n^{(1)}_{uv},\; n^{(2)}_{uv},\; x_u - x_v\,] \in \mathbb{R}^{3\times 3}.
\]}

The length was a scalar sharing a tensor with directions, which is exactly the
violation described in \S\ref{sub:requirements}. Its content --- local scale --- is recovered by
the relative position $x_u - x_v$, whose norm \emph{is} the length and which
additionally carries direction. The midpoint duplicated information already
present in the two endpoint coordinates.

The two face normals are stored in a canonical order derived from the mesh
orientation, so that the pair is a well-defined function of the edge rather than
of the order in which faces happened to be enumerated. Only these two channels
participate in the face-normal loss; the relative position must not, because a
cosine loss would normalise away the very length that is its content.

\subsection{Fragment scale as an invariant gate}
\label{sub:scale}

Fragments span four to 83{,}039 vertices, so raw scale is heavy-tailed and would
dominate anything it were concatenated with. Scale has no direction, so it
cannot join the vector stream. It conditions that stream multiplicatively
instead:
\[
g = \operatorname{softplus}\!\big(\mathrm{MLP}([\,\mathrm{Gram}(x),\ \log s\,])\big),
\qquad x \leftarrow g \odot x ,
\]
where $\mathrm{Gram}(x)$ denotes the matrix of pairwise inner products of the
vector channels, which is invariant, and $s$ is the fragment scale. The gain $g$
is invariant, so the product remains equivariant. The logarithm is used because
of the heavy tail; GARF applies a positional encoding for the same reason.

\subsection{Connectivity}
\label{sub:connectivity}

\change{V4 retained both an adjacency (COO) structure and an incidence
structure pairing nodes with globally shifted edge identifiers. Only the
adjacency structure is retained.}

Edges are stored bidirectionally, so structural symmetry is preserved. The
incidence structure existed to let messages pass between nodes and edges; in the
implemented design, edge features are consumed directly by the attention
mechanism at each edge, which serves the same purpose without a second index
tensor to keep synchronised.

A correctness requirement that is easy to miss: the ground-truth and perturbed
copies of a scene must share \emph{one} topology. The losses match vertices,
edges and normals by index, so two independently constructed copies would
compare unrelated rows without raising an error. Face de-duplication reorders
faces lexicographically and can produce two different vertex orderings, so the
two copies are built once and split, and their edge indices are asserted equal.

\section{Equivariance: the constraint that shapes the design}
\label{sec:equivariance}

\subsection{Definition}
\label{sub:definition}

Let $R \in \SOthree$. A mapping $L$ is equivariant when
\[
L(R \cdot f_k) = D(R)\, L(f_k),
\]
where $D(R)$ acts on the output space. A conventional network treats two rotated
copies of one fragment as unrelated inputs and must learn the same geometry many
times over; an equivariant layer obtains the rotation for free, as an algebraic
property rather than as something acquired from augmented data.

\subsection{Equivariant to its own pose, invariant to every other fragment's}
\label{sub:twolines}

\change{V4 stated only the equation above. It is necessary but not sufficient,
and the missing half determines the entire inter-fragment design.}

Each fragment is perturbed by its \emph{own} rotation. The label for fragment
$i$ therefore does not change when fragment $j$ lands at a different angle --- and
the prediction must not either. Writing $A$ for a rotation applied to a single
fragment's input, fragment $i$'s representation $H_i$ must satisfy both lines:
\[
\begin{aligned}
H_i(P_1,\ldots,P_i A,\ldots,P_N) &= H_i(\,\cdot\,)\, A &&\text{equivariant to its own pose,}\\
H_i(P_1,\ldots,P_j A,\ldots,P_N) &= H_i(\,\cdot\,) &&\text{invariant to every other pose } (j \neq i).
\end{aligned}
\]

This is equation 7 of Wu et al., \emph{Leveraging SE(3) Equivariance for
Learning 3D Geometric Shape Assembly}. The second line is forced by the task,
before any architecture is chosen. An architecture that can see fragment $j$'s
orientation must \emph{learn} to ignore it --- from data, by augmentation ---
which is precisely what equivariance is supposed to make unnecessary. Worse,
$j$'s orientation is not weak signal but pure noise injected by the
perturbation, feeding directly into $i$'s prediction.

\subsection{What may cross the gap}
\label{sub:crosses}

The constraint is not ``no vectors between fragments''. The cross-fragment layer
takes $(T,C,3)$ vector features and returns $(T,C,3)$ vector features, and every
operation on the spatial axis is equivariant. What is constrained is what a
message may \emph{depend on}.

Wu et al.'s correlation module is the construction that satisfies both lines:
$C_{ij} = G_j \cdot F_i$, where $G_j$ is an \emph{invariant matrix} describing
fragment $j$ and $F_i$ is fragment $i$'s \emph{own equivariant} features. A
matrix, not a scalar --- the sender may re-mix the receiver's channels
arbitrarily. What it cannot do is contribute a direction of its own, because a
direction carries an orientation.

What crosses is \emph{shape}. What does not cross is \emph{pose}. Nothing is
lost by that: the relative pose of two arbitrarily tumbled fragments is the
perturbation, not the object.

A corollary constrains the attention scores. The natural score
$\langle W_q x_i, W_k x_j\rangle$ is unusable, because under independent
perturbations it becomes $\langle A_i a, A_j b\rangle$, which depends on the
relative rotation $A_i^{\top}A_j$ --- exactly the quantity that is noise. Scores
are therefore inner products of each token's invariant description.

\subsection{A consequence: the perturbation cancels from the error}
\label{sub:cancels}

This property has a consequence for reading training curves that is worth
stating explicitly, because the obvious interpretation of an initial loss is
wrong.

The head returns $R_{\text{pred}} = M^{\top}$, where $M$ is the equivariant
frame, and the frame co-rotates with its input. On a fragment perturbed by $Q$
the prediction is $M_{\text{asm}}^{\top} Q^{\top}$ while the label is
$Q^{\top}$, so the error rotation is
\[
R_{\text{pred}} R_{\text{label}}^{\top}
= M_{\text{asm}}^{\top} Q^{\top} Q
= M_{\text{asm}}^{\top}.
\]

\textbf{The perturbation cancels exactly.} An untrained model's rotation error is
the geodesic angle of its own frame on the \emph{assembled} fragment, and has
nothing to do with how that fragment was tumbled. Verified numerically: across
three independent perturbations the per-fragment errors agreed to
\measured{$3\times10^{-13}$ radians}.

Two things follow. First, $126.5^{\circ}$ is not a distribution an untrained
equivariant model draws from; it is the mean angle between two \emph{independent}
uniform rotations. Here the error is a deterministic function of the untrained
network, so it equals $126.5^{\circ}$ only if those frames happen to be
uniformly spread. Measured over twelve seeds on a 24-fragment batch, the initial
error was $132.3^{\circ}\pm11.8$ for one architecture and $121.9^{\circ}\pm15.0$
for another --- \emph{architecture-dependent}, not a defect.

Second, it restates the objective usefully. Since the initial error \emph{is}
$\angle M_{\text{asm}}$, training is literally driving the frame on assembled
fragments to the identity. That is a fixed, learnable target, because the
assembled pose is a convention shared by every scene in the dataset.

\subsection{Numerical verification}
\label{sub:verification}

Equivariance is an algebraic identity, so it should hold to machine precision
rather than approximately. All checks below were run in double precision.

\begin{table}[h]
\centering
\small
\begin{tabular}{lr}
\toprule
\textbf{Property} & \textbf{Residual} \\
\midrule
Every Vector Neuron primitive, $L(xR^{\top}) = L(x)R^{\top}$ & $\leq 10^{-15}$ \\
Gram--Schmidt head orthonormal, $\det = +1$ & exact to $10^{-16}$ \\
Intra-fragment attention, equivariance & $1.3\times10^{-15}$ \\
Cross-fragment, own-pose equivariance & $4.4\times10^{-16}$ \\
Cross-fragment, other-pose \emph{invariance} & $3.3\times10^{-16}$ \\
Cross-scene leakage within a batch & exactly $0$ \\
Whole model, per-fragment equivariance & $\leq 10^{-9}$ \\
Perturbation cancellation (\S\ref{sub:cancels}) & $3\times10^{-13}$ \\
\bottomrule
\end{tabular}
\caption{Measured equivariance residuals. A failure at $10^{-6}$ would indicate
a real break rather than accumulated round-off.}
\end{table}

Two defects were caught by these checks, both of the kind that trains without
complaining. A near-identity initialisation implemented by zeroing a gate's
final weight makes the layer's output constant, which zeroes the gradient into
everything upstream --- 31 of 77 parameters dead on the first steps, waking only
as a bias drifted. And $\sqrt{s + \epsilon}$ perturbs \emph{every} norm rather
than only the near-zero ones: the head came out orthonormal to only $10^{-8}$
with $\det = 0.99999995$. Clamping the squared sum instead is exact away from
zero.

\section{Fracture surfaces and inter-fragment tokens}
\label{sec:tokens}

\subsection{Why the virtual-node design was withdrawn}
\label{sub:virtual}

V4 introduced eight coordinate-aware virtual nodes per fragment, pooled upward
from vertices, attended across fragments, and broadcast back down. That design
is withdrawn for two reasons.

The first is the constraint of \S\ref{sub:twolines}. Virtual nodes with fixed spatial
coordinates are \emph{positions}, and pooling vertex features onto them and
attending between them lets a fragment's orientation influence its neighbours'
representations. The mechanism cannot satisfy the invariance line without
additional machinery.

The second is empirical. The earlier VN-GAT design reached approximately
$43^{\circ}$ on training data but stalled near $94^{\circ}$ on validation --- close
to the \emph{axis-only floor} of $89.9^{\circ}$, which is the score of a model
that recovers a fragment's axis but not its rotation about that axis. That is
the signature of a per-fragment canonicaliser with no usable access to
inter-fragment relationships.

\change{Virtual nodes are replaced by direct attention between vertices lying on
fracture surfaces. All vertices are retained as graph nodes; the fracture mask
decides only which vertices are \emph{eligible} to become cross-fragment
tokens.}

\subsection{Two labellings, and the inference gap}
\label{sub:labellings}

A fragment's surface is two things glued together: the object's original
exterior, which is smooth and shared with no one, and the fracture surface,
where it tore away and which matches exactly one other fragment. Two independent
labellings are available.

\begin{table}[h]
\centering
\small
\begin{tabular}{lll}
\toprule
\textbf{Method} & \textbf{What it uses} & \textbf{Available at inference?} \\
\midrule
Dihedral & sharp dihedral variation within one fragment & yes, one fragment alone \\
Coincidence & vertices another fragment also occupies & no, needs the assembled scene \\
\bottomrule
\end{tabular}
\end{table}

Coincidence is ground truth and is how supervision is derived. Dihedral is the
estimate a model can compute at test time, and it is \measured{much weaker than
V4 assumed}: at its tuned threshold it over-labels by approximately
\measured{3.9$\times$}, with precision 0.24 against recall 0.93. It also lights
up handles, rims and feet --- sharp because the object was \emph{designed} that
way, not because it broke there.

This is the single largest gap between the reported results and a real
deployment, and \S\ref{sec:status} states it as such. The threshold is chosen by measurement
rather than by eye, scoring dihedral against coincidence across a sweep and
reporting recall-weighted $F_\beta$ rather than $F_1$: a missed fracture face is
an edge the model never sees, while a spurious one is only noise, so the two
error types are not symmetric.

\subsection{Token selection: geodesic-disk sampling under a per-scene budget}
\label{sub:sampling}

Cross-fragment tokens are selected from the masked vertices by farthest-point
sampling under a \emph{geodesic} metric --- shortest path along fracture-surface
mesh edges --- rather than a Euclidean one. The distinction matters on folded or
re-entrant surfaces, where two points may be close in space while lying far
apart on the surface; Euclidean sampling treats them as redundant and leaves
part of the interface unsampled.

The budget is \emph{per scene}, not per fragment. A single total $T$ is divided
among the fragments in proportion to their fracture-surface area, and sampling
then proceeds within each fragment's own eligible set. Two properties follow.

A per-fragment cap is deliberately \emph{not} imposed. If a scene contains one
large fragment and several small ones --- a head with the ears, nose and part of
the hair broken off --- proportional allocation gives the large piece more
tokens, which is the intended behaviour. Capping would destroy that
proportionality for no benefit, because the total alone already bounds the cost:
for any split $\{t_i\}$ summing to $T$, the number of undirected cross-fragment
pairs is
\[
\frac{T^2 - \sum_i t_i^2}{2} \;<\; \frac{T^2}{2}.
\]

Sampling is per fragment and never over the pooled scene. At input time the
fragments sit in arbitrary perturbed poses, so a single global farthest-point
pass would select points according to where the pieces happened to land --- a
different sample for every perturbation of the same scene, and no equivariance
left.

\subsection{The cost of the resulting graph}
\label{sub:graphcost}

V4's virtual-node compression was motivated by an assumed cost. With the true
fracture surface measured at 10.6\% of vertices, connecting every fracture
vertex to every fracture vertex of every \emph{other} fragment turns out to be
affordable with no reduction at all.

\begin{table}[h]
\centering
\small
\begin{tabular}{lrr}
\toprule
\textbf{Fragments in scene} & \textbf{Cross-fragment pairs} & \textbf{vs GARF's 6-layer stack} \\
\midrule
3 (median scene) & 46{,}875 & 0.001$\times$ \\
8 (mean scene) & 753{,}088 & 0.01$\times$ \\
20 & 2{,}968{,}750 & 0.04$\times$ \\
99 (largest in corpus) & 75{,}796{,}875 & 1.01$\times$ \\
\bottomrule
\end{tabular}
\caption{Unreduced inter-fragment graph cost. It crosses GARF's whole stack only
at 99 fragments, which is exactly the largest scene in the dataset.}
\end{table}

\textbf{Reduction is therefore not needed for cost.} Sampling is retained for a
different reason: to make the inter-fragment connection count \emph{fixed and
predictable} rather than a function of whatever mesh happens to arrive. A
retraction is owed here, and is recorded in \S\ref{sec:status}.

\section{Network architecture}
\label{sec:architecture}

\subsection{Vector Neuron primitives}
\label{sub:vn}

Features are $(C,3)$ tensors. The primitives are bias-free linear channel
mixing, a learned-hyperplane leaky rectifier, a layer normalisation over
invariant magnitudes, and an invariant readout formed from inner products with
learned directions. A linear map with a bias would break equivariance, because a
constant offset does not rotate; the absence of bias is structural, not
stylistic.

\subsection{Intra-fragment attention}
\label{sub:intra}

Attention along real mesh edges. Scores are invariant inner products of
projected query and key features, so weighting equivariant messages by them
leaves the result equivariant. Each node also carries a self-transform outside
the softmax, so a node with no incident edges still has an output, and so the
layer can learn to ignore its neighbourhood entirely.

Two implementation properties are worth recording because they are not
cosmetic. Because the linear maps have no bias,
$W[a;b;c] = W_a a + W_b b + W_c c$ exactly, so concatenation can be replaced by
separate projections that are summed. This removes the concatenation temporaries
and, more importantly, allows each projection to run on the $N$ \emph{nodes}
before being gathered to the $E$ \emph{edges} rather than after --- roughly one
sixth of the work on a triangle mesh.

\subsection{Cross-fragment attention}
\label{sub:cross}

The layer that replaces V4's virtual nodes. It is a Vector Neuron layer
throughout, and it implements the correlation construction of \S\ref{sub:crosses}:
\[
\begin{aligned}
\text{logits}_{ij} &= \langle q(\mathrm{inv}_i),\, k(\mathrm{inv}_j)\rangle / \sqrt{d}
&&\text{invariant to both poses,}\\
\alpha &= \operatorname{softmax}_j \text{ within the scene, excluding } j\text{'s own fragment,}\\
\mathrm{out}_i &= x_i + \Big(\sum_j \alpha_{ij} G_j\Big) x_i .
\end{aligned}
\]
$G_j$ is a per-head channel-mixing matrix read from fragment $j$'s invariants.
Because $G$ is a \emph{linear} function of those invariants,
$\sum_j \alpha_{ij} G(v_j) = G\big(\sum_j \alpha_{ij} v_j\big)$, so the
aggregation happens on the small invariant vectors and the matrix is built once
per query rather than once per pair. This is not a cosmetic rearrangement: a
per-pair matrix would be gigabytes at the pair counts a large scene reaches.

Tokens attend across fragments \emph{within one scene}, never inside their own
fragment --- the mesh graph already covers that --- and never across scenes. The
cross-scene case is not hypothetical: a batch concatenates unrelated objects and
fragment identifiers are globally unique, so an unmasked layer silently lets one
object inform another and makes the prediction depend on batch composition. That
defect occurred during development and is now excluded by construction and
asserted in the test suite.

\subsection{Layer schedule and propagation reach}
\label{sub:schedule}

The backbone is five intra-fragment layers, three cross-fragment layers, then
one further intra-fragment layer: describe each fragment, let the fragments
talk, then spread what they heard. The final clause is not decoration.

A cross layer writes only to \emph{token} vertices, and the rotation head pools
a mean over \emph{every} vertex --- so a vertex the tokens never reach
contributes to the prediction without having heard from another fragment. Each
intra layer after the last cross layer buys exactly one hop along mesh edges.
Measured on real scenes:

\begin{table}[h]
\centering
\small
\begin{tabular}{lr}
\toprule
\textbf{After the last cross layer} & \textbf{Vertices informed} \\
\midrule
nothing & 17\% \\
one intra layer & \measured{49\%} \\
two intra layers & 64\% \\
\bottomrule
\end{tabular}
\caption{Fraction of vertices within reach of a cross-fragment token. The first
trailing layer buys 32 points, the second only 15 --- tokens cluster on fracture
surfaces, so their neighbourhoods overlap.}
\end{table}

The reach is measured on the actual meshes at start-up rather than assumed,
because it depends on how the fracture surface is shaped and cannot be derived
from the token count alone.

\subsection{Rotation head}
\label{sub:head}

\change{V4 predicted a normalised quaternion. V5 predicts two equivariant
3-vectors and orthonormalises them by Gram--Schmidt.}

The quaternion was chosen in V4 to avoid gimbal lock. It does, but it shares
with every three- and four-parameter parameterisation a subtler problem: the map
from $\SOthree$ to the representation is discontinuous, so a network must learn
a discontinuous function, and the antipodal identification $q \sim -q$ leaves
the target ambiguous unless handled explicitly. The 6D representation is
continuous, which is the property that matters for regression.

More decisively, it fits the equivariant setting exactly: the two vectors are
ordinary vector features, so they rotate with the input for free, and
Gram--Schmidt is equivariant. The head therefore produces an equivariant frame
$M$ with no additional constraint, and returns $R_{\text{pred}} = M^{\top}$
(the \emph{rows}, not the columns). The derivation, and the reason this
transpose is the single easiest thing in the system to get backwards, are given
in \S\ref{sub:rotloss}.

\section{Training objective}
\label{sec:objective}

\subsection{Rotation loss}
\label{sub:rotloss}

Centred, the perturbation satisfies $v_{\text{pert}} = v_{\text{gt}} Q^{\top}$.
Vector features inherit that rotation, so the frame satisfies
$M_{\text{pert}} = Q M_{\text{gt}}$. The label is the rotation that undoes the
perturbation, $R_{\text{gt}} = Q^{\top}$, so the head returns
$R_{\text{pred}} = M^{\top}$, giving
$R_{\text{pred}} = M_{\text{gt}}^{\top} Q^{\top}$ --- correct exactly when the
network maps an assembled fragment to the identity frame.

Getting this transpose backwards is invisible at initialisation, because chance
is chance in either direction, and would surface only as a model that never
converges. It is therefore asserted directly against a known perturbation rather
than inferred from a loss curve.

\change{V4 defined the angular error as
$\theta = \arccos\big((\operatorname{tr}(R_k^{\top}R_{\text{gt}})-1)/2\big)$.
V5 uses an $\operatorname{atan2}$ formulation of the same angle.}

The two agree in value and differ in gradient. The derivative of $\arccos$ is
unbounded at its endpoints, so the gradient of V4's loss \emph{diverges exactly
as the model converges}, swamping every other term at precisely the point where
fine alignment matters. The $\operatorname{atan2}$ form, taking the angle from
the norm of the axis component against the trace, is finite throughout. Its
guarded norm also takes a far smaller floor, because $\operatorname{atan2}$'s
derivative cancels the norm's --- the loss is exactly zero at a perfect fit
rather than reaching a $5\times10^{-9}$ radian floor.

\subsection{Geometric auxiliary losses}
\label{sub:aux}

The predicted rotation is applied to the centred, non-standardised fragment and
compared against the ground-truth copy, index by index:
\[
\begin{aligned}
L_{\text{pos}} &= \frac{1}{|V|}\sum_{v} \lVert x_v - \hat{x}_v \rVert_2^2,
&\qquad
L_{\text{normal}} &= \frac{1}{|V|}\sum_{v} (1 - n_v \cdot \hat{n}_v), \\
L_{\text{face}} &= \frac{1}{|E|}\sum_{e} \Big[(1 - n^{(1)}_e \cdot \hat{n}^{(1)}_e) + (1 - n^{(2)}_e \cdot \hat{n}^{(2)}_e)\Big] .
\end{aligned}
\]

\change{V4's edge-midpoint loss $L_{\text{mid}}$ is removed together with the
midpoint feature.}

The face term is restricted to the two normal channels inside the loss rather
than by the caller. Edges carry three channels, and passing all three is the
obvious mistake: a cosine loss would normalise the relative position, whose
\emph{length} is its content, and would silently move the term's chance value
from 2.0 to 3.0.

\subsection{Interface correspondence: two corrections}
\label{sub:corr}

This is the term that has changed most, and both changes are corrections.

\textbf{V4's formulation is minimised by cancellation.} V4 defined
\[
L_{\text{emb-}v} = \sum_{x \in X_{\text{shared}}}
\Big\lVert \sum_{(f,v) \in C_V(x)} z^{(f)}_v \Big\rVert_2^2 .
\]
Two \emph{opposite} embeddings score zero; two \emph{identical} ones score
$4\lVert z\rVert^2$. The objective rewards exactly the wrong configuration while
looking like a reasonable consistency penalty.

\textbf{The obvious repair is minimised by collapse.} Replacing the norm of the
sum with the scatter about each cluster centroid,
\[
L = \operatorname{mean}_c \operatorname{mean}_{i \in c}\, \lVert z_i - \bar{z}_c \rVert^2 ,
\]
fixes the cancellation and leaves a second trivial minimum: a \emph{constant}
embedding scores exactly zero. This is not hypothetical --- it is what the first
real training run did. Over 180 steps the loss fell from 0.0109 to 0.0010 while
the spread of the embeddings fell from 0.166 to 0.021 and their norm held at
2.0. The head was converging on a single vector and reporting success the whole
way down. A constant embedding is worse than an untrained one for what this head
exists to do, since stage two matches by mutual nearest neighbours and a
constant embedding ties every distance.

\change{The term is now InfoNCE over the coincidence clusters:
\[
L_{\text{corr}} = -\operatorname{mean}_i \log
\frac{\sum_{p \in \mathrm{pos}(i)} \exp(s_{ip}/\tau)}
     {\sum_{k \neq i} \exp(s_{ik}/\tau)},
\]
with $s$ the cosine similarity between $\ell_2$-normalised embeddings.}

It is written contrastively rather than as a variance floor because the
negatives are exactly the confusion set at matching time --- the other fracture
vertices of the same scene --- so the training objective and the downstream use
become the same question. Collapse now scores $\log(A-1) - \log|\mathrm{pos}|$,
verified against the closed form: \measured{4.844} for 128 anchors in pairs,
against \measured{0.024} for a perfect embedding. The point is not that collapse
became impossible, but that it became \emph{visible}.

Anchors are capped per batch, because the similarity matrix is quadratic in them
and a scene can label thousands of vertices.

A reporting consequence: the composite total no longer bottoms out at zero,
since InfoNCE retains a small temperature-dependent floor even for a perfect
embedding. The number to read instead is \textbf{match@1}, the fraction of
coincident vertices whose true partner is their own nearest neighbour. It runs
from near zero at chance to exactly 1, needs no per-batch reference, and a
collapsed embedding cannot fake it.

\subsection{Composite objective}
\label{sub:composite}

\[
L_{\text{total}} = \lambda_1 L_{\text{rot}} + \lambda_2 L_{\text{pos}}
+ \lambda_3 L_{\text{normal}} + \lambda_4 L_{\text{face}}
+ \lambda_5 L_{\text{corr}} .
\]

All weights are currently 1.0 and \emph{deliberately untuned}. The terms have
very different natural scales --- radians, a cosine in $[0,2]$, a distance in
normalised units, a log-ratio --- and choosing weights before observing how each
one moves is guesswork presented as a decision. Every term is reported
separately so that the first training runs \emph{measure} the relative scales
rather than assume them.

\section{Reference values}
\label{sec:reference}

\change{V4 specified no reference values. Every number in this system is now
read against what it must be when the model has learned nothing.}

A geometric loss that is silently wrong still descends, so ``the loss went
down'' proves nothing. All values below were established by Monte Carlo with
fixed seeds, not asserted.

\begin{table}[h]
\centering
\small
\begin{tabular}{llp{7.2cm}}
\toprule
\textbf{Quantity} & \textbf{At chance} & \textbf{Meaning} \\
\midrule
Geodesic angle & $126.48^{\circ}$ & $\pi/2 + 2/\pi$; a model here has learned nothing \\
Euler RMSE, random & $86.29^{\circ}$ & an independent uniform prediction \\
Euler RMSE, identity & $83.14^{\circ}$ & always predicting no rotation \\
Axis-only floor & $89.9^{\circ}$ & axis recovered, rotation about it not \\
Node normal loss & $1.0$ & mean of $1 - \cos$ over random directions \\
Face normal loss & $2.0$ & two normals per edge, hence twice \\
Node position loss & $4/3$ & on a unit sphere \\
\bottomrule
\end{tabular}
\caption{Reference values. Every term of the objective has one.}
\end{table}

\textbf{The two Euler rows are the important pair, and the ambiguity in the word
``chance'' is not harmless.} A model that collapses to predicting the identity
--- the cheapest way to reduce a rotation loss early in training --- scores
$83.14^{\circ}$, which is \emph{three degrees better than guessing}, on the
metric GARF reports as its headline. The geodesic angle is identical for both,
$126.48^{\circ}$, because a uniform rotation's angle is distributed the same way
whether measured against the identity or against another uniform rotation. The
geodesic angle therefore does not pay for the collapse, which is the argument
for reporting it as the primary number and Euler RMSE only for comparability.

The initial rotation reading is checked against a deliberately wide band, for
the reason established in \S\ref{sub:cancels}: it is the untrained frame's own angle and
moves with the architecture. A band tight around $126.5^{\circ}$ would flag
every architecture change as a defect. What the check still catches is gross
breakage --- a model beginning near $0^{\circ}$ or $180^{\circ}$ cannot be
measuring the angle it claims to.

\section{Correspondence discovery}
\label{sec:correspondence}

After training, interface embeddings identify candidate matching points between
fragments. For every interface point a nearest-neighbour search is performed in
embedding space, and \emph{mutual} nearest-neighbour matching is required: two
points are a correspondence only if each selects the other. This reduces false
matches and provides reliable geometric constraints for the translation solver.
This section is unchanged from V4 in principle; what has changed is that the
embedding it operates on is now trained by an objective that cannot be satisfied
by a constant (\S\ref{sub:corr}).

\section{Geometric optimisation-based translation recovery}
\label{sec:translation}

Translation estimation differs fundamentally from rotation estimation.
Determining orientation requires semantic understanding of geometry and
fragment relationships; determining translation, once orientations are known, is
primarily a geometric consistency problem. Translations are therefore not
predicted by the network but obtained by solving an optimisation.

Let $p_i$ and $q_i$ denote corresponding interface points from mutual
nearest-neighbour matching.

\begin{align*}
E_{\text{pos}} &= \sum_i \lVert p_i - q_i \rVert^2
&&\text{matching surfaces should occupy the same locations,}\\
E_{\text{normal}} &= \sum_i \lVert n^{(A)}_i + n^{(B)}_i \rVert^2
&&\text{mating surfaces should have opposing normals,}\\
E_{\text{collision}} &&&\text{penalises volume interpenetration.}
\end{align*}

The translation vectors are recovered globally by
\[
\min_{t_1,\ldots,t_N}\ \Big( E_{\text{pos}} + \lambda_{E_1} E_{\text{normal}}
+ \lambda_{E_2} E_{\text{collision}} \Big),
\]
placing fragments so as to satisfy interface alignment, normal consistency and
collision avoidance simultaneously.

\textbf{This stage is specified but not implemented.} It is stated here as
unchanged from V4 so that the design remains complete, but no result in this
document depends on it, and \S\ref{sec:status} records it as outstanding.

\section{Workflow: training, validation, and inference}
\label{sec:workflow}

\subsection{The shared path}

Training and validation share almost the entire pipeline. A scene directory is
read once into fragments in their assembled frame; the fracture mask is computed
from geometry; coincidence clusters and the ground-truth copy are taken in the
assembled frame; a random rotation is drawn per fragment; the scene is
standardised and tokens are sampled; fragments are concatenated into one ragged
batch. The network runs identically in both modes. The path diverges only at the
end: training computes the objective and steps the optimiser, while validation
runs the same forward pass without gradients and scores the prediction.

Two ordering constraints are absolute. \textbf{Coincidence must be computed
before perturbation}, because it is defined in the assembled frame; afterwards
nothing coincides and the labels come back empty while looking like a result.
And the ground-truth and perturbed copies must be built once and split, for the
topology reason of \S\ref{sub:connectivity}.

\subsection{Validation is not label-free}

The perturbation is generated by the pipeline itself, so $Q^{\top}$ is known and
every metric is computable. What validation measures is generalisation to unseen
\emph{objects}, not to unknown labels.

A defect found here is worth recording, because it is the worst shape a defect
can take. Official split matching originally used only the object directory
name, discarding the category. Breaking Bad lists objects as
\texttt{category/object} and directory names are \emph{not} unique across
categories, so sixteen objects appeared in both the training and validation
splits. Validation then becomes partly a memorisation test, and the number it
produces is \emph{better} than the honest one --- it reads as success. Matching
is now on the pair, and the loader refuses to run if any two splits intersect.

\subsection{A real broken object}

A third regime is worth separating explicitly, because it is the one the thesis
must be honest about. On a genuinely broken object rather than a Breaking Bad
sample:

\begin{itemize}
\item there is no $Q$, hence no label --- a prediction exists and nothing checks it;
\item there are no coincidence labels, so the dihedral mask alone selects tokens,
      at precision 0.24 (\S\ref{sub:labellings}). GARF's answer is a trained fracture
      segmenter reaching above 99.5\%; this project does not have one;
\item a rotation does not place a fragment --- stage two is required and is not built.
\end{itemize}

Any figure reported on Breaking Bad is therefore an \emph{optimistic bound},
because it uses ground-truth fracture labels a deployment would not have.

\section{Implementation within a commodity hardware budget}
\label{sec:implementation}

The economy question of \S\ref{sec:intro} makes implementation a design concern rather than
an afterthought, so the constraints that shaped it are recorded here.

\textbf{The binding constraint is memory, and it was in an unexpected place.}
Cross-fragment attention gathers query, key and value tensors to the \emph{pair}
dimension. Measured as the slope against pair count, that retained
\measured{1{,}109 bytes per pair}: at 2{,}048 tokens over six fragments ---
roughly 3.5 million pairs --- \measured{3.6 GB for a single scene}, per layer.
Recomputing those gathers during the backward pass instead of storing them costs
one extra forward pass of an indexing operation and takes the figure to
\measured{86 bytes per pair, 0.28 GB}, a thirteen-fold reduction with outputs
and gradients bitwise identical.

\textbf{A retraction belongs here.} The token budget of 2{,}048 per scene was
originally justified by comparing pair counts against GARF's stack --- a
\emph{floating-point} argument, made about hardware with far more memory. It was
never checked against the memory of the target device, which is the binding
constraint and is quadratic in exactly that number. The value survives, but it
survived by luck rather than because the check had been performed.

\textbf{Gradient accumulation is not a free substitute for a larger batch.} When
memory forces the batch down, accumulation is the standard remedy, but dividing
each micro-batch's loss by the accumulation count weights each \emph{scene}
equally, while a true batch weights each \emph{fragment} equally --- and a scene
holds between two and 35 fragments. Measured on a two-fragment and an
eight-fragment scene, the two gradients had cosine similarity \measured{0.80}
and norms 45\% apart: a silent re-weighting of the objective introduced by a
setting chosen for memory reasons. Weighting each micro-batch by its fragment
count and normalising by the group total makes the two identical to seven
decimal places on all four geometric terms. The contrastive term cannot be made
to match, because its negatives are drawn from the batch --- and should not be,
since within-scene negatives are the confusion set that matching actually faces.

\textbf{Sessions are capped, so a run is a chain.} Training must survive being
killed at a fixed wall-clock limit: checkpoints are written on a timer rather
than only at epoch boundaries, because a single epoch over the full training set
can outlast a session; the always-current checkpoint is kept separate from the
best-so-far, because a policy-gated file alone can freeze while training
continues and a resume then silently discards the gap; and the learning-rate
schedule is a pure function of the step count rather than a stateful object,
because such objects are periodic past their horizon and store that horizon in
their own state.

\section{Status, and what has not been established}
\label{sec:status}

This section exists because the failure mode of this class of system is a
plausible number rather than an exception.

\subsection{Established}

The data and geometry layer, the equivariant network, the composite objective
and the training engine are implemented and covered by an automated test suite
that runs clean under warnings-as-errors. Equivariance holds to the residuals of
\S\ref{sub:verification}. The rotation convention is asserted by round trip rather than inferred.
The architecture demonstrably \emph{can} fit: driven on a small fixed set it
reduces rotation error from $113^{\circ}$ to \measured{$0.4^{\circ}$} on one
scene and from $137^{\circ}$ to \measured{$0.6^{\circ}$} on three; on 32 scenes
it reached $13.3^{\circ}$ after 1{,}608 steps and was still descending. These
are \emph{training}-error results and prove nothing about generalisation --- which
is exactly what makes them useful, since they isolate ``are the conventions
self-consistent'' from ``does it generalise''.

\subsection{Not established}

\textbf{Nothing has been trained to convergence on the real dataset.} Every
figure in this document is a property of the architecture, the data, or the
implementation --- not evidence that the system learns the task. A short
calibration run over two full epochs left the rotation error flat at chance,
which at roughly 3{,}500 of a planned 65{,}000 optimiser steps is not yet
evidence either way.

\subsection{Open questions}

\begin{itemize}
\item \textbf{Inference-time labelling.} Coincidence is unavailable at test
      time and dihedral is weak (\S\ref{sub:labellings}). Until a segmenter exists, reported
      numbers are an optimistic bound.
\item \textbf{Supervision asymmetry.} The correspondence loss is defined over
      coincident vertices, so it supervises only the fracture channel. Tokens
      drawn from the original surface are shaped by the rotation loss alone.
\item \textbf{Loss weights, layer counts, token budget.} All are currently
      chosen rather than measured. The weights are deliberately equal so that
      the first runs measure their natural scales.
\item \textbf{Fragments with no fracture surface.} Coincidence reports a minimum
      of zero, which should be impossible; most likely single-fragment modes,
      but it needs counting, because such fragments receive no inter-fragment
      edges at all.
\item \textbf{The translation solver.} Specified in \S\ref{sec:translation}, not implemented.
\end{itemize}

\end{document}
